\documentclass[12pt]{article}
\usepackage{amsmath,amssymb}
\usepackage{geometry}
\usepackage{enumitem}
\geometry{margin=0.8in}

\title{Algebra II Review: Complex Numbers as Plane Vectors}
\author{}
\date{}

\begin{document}
\maketitle

\section*{Context}
A complex number behaves like a two-dimensional real vector, but it also has a specially defined multiplication. That multiplication turns each pair $\langle a,b\rangle$ into a single number that can be added, multiplied, divided, measured, and assigned a direction. This viewpoint explains the usual notation $a+bi$ rather than treating it as an arbitrary rule.

\section*{Real Pairs with Special Operations}
Begin with ordered pairs of real numbers:
\[
  z=\langle a,b\rangle,
  \qquad
  w=\langle c,d\rangle.
\]
Addition is ordinary vector addition:
\[
  \langle a,b\rangle+\langle c,d\rangle
  =\langle a+c,b+d\rangle.
\]
Define multiplication by the special rule
\[
  \boxed{\langle a,b\rangle\langle c,d\rangle
  =\langle ac-bd,ad+bc\rangle.}
\]
This is \emph{not} component-by-component multiplication.

If $r$ is a real scalar, identify it with the pair $\langle r,0\rangle$. Multiplication by this pair gives ordinary scalar multiplication:
\[
\begin{aligned}
  \langle r,0\rangle\langle a,b\rangle
  &=\langle ra-0b,rb+0a\rangle\\
  &=\langle ra,rb\rangle
  =r\langle a,b\rangle.
\end{aligned}
\]
Thus real scalar multiplication is a special case of complex multiplication.

\section*{Length and Direction}
The length, or \textbf{modulus}, of $z=\langle a,b\rangle$ is its distance from the origin:
\[
  |z|=\sqrt{a^2+b^2}.
\]
The \textbf{argument} of a nonzero complex number is a direction angle $\theta$ measured from the positive real axis:
\[
  \arg z=\theta,
  \qquad
  \cos\theta=\frac{a}{|z|},
  \qquad
  \sin\theta=\frac{b}{|z|}.
\]
The quadrant must agree with the signs of $a$ and $b$. Arguments differing by $2\pi k$ describe the same direction. In this sheet, $\operatorname{Arg}z$ denotes the principal argument in $0\le\theta<2\pi$.

If $r=|z|$, then
\[
  z=\langle r\cos\theta,r\sin\theta\rangle.
\]

\section*{From $\langle a,b\rangle$ to $a+bi$}
Define
\[
  1:=\langle1,0\rangle,
  \qquad
  i:=\langle0,1\rangle.
\]
Every pair can now be written as
\[
\begin{aligned}
  \langle a,b\rangle
  &=a\langle1,0\rangle+b\langle0,1\rangle\\
  &=a\cdot1+b\cdot i\\
  &=a+bi.
\end{aligned}
\]
Therefore
\[
  \boxed{\langle a,b\rangle\longleftrightarrow a+bi.}
\]
The pair and $a+bi$ are two representations of the same object. Geometrically, $a+bi$ is plotted at the point $(a,b)$ in the complex plane. Algebraically, it is one scalar complex number. Complex scalars are therefore ``like'' two-dimensional real vectors, but their special multiplication gives them additional algebraic structure.

The multiplication rule explains why $i$ is a square root of $-1$:
\[
\begin{aligned}
  i^2
  &=\langle0,1\rangle\langle0,1\rangle\\
  &=\langle0\cdot0-1\cdot1,\,0\cdot1+1\cdot0\rangle\\
  &=\langle-1,0\rangle=-1.
\end{aligned}
\]
Thus $i^2=-1$, so $i$ is a number whose square is $-1$. In the usual Algebra II notation, $i=\sqrt{-1}$.

Using $i^2=-1$, pair multiplication becomes the familiar calculation
\[
  (a+bi)(c+di)=(ac-bd)+(ad+bc)i.
\]

\section*{Real and Imaginary Parts}
For $z=a+bi$,
\[
  \operatorname{Re}(z)=a,
  \qquad
  \operatorname{Im}(z)=b.
\]
The imaginary part is the real coefficient $b$, not the term $bi$. The pair notation records the same information:
\[
  z=\langle\operatorname{Re}(z),\operatorname{Im}(z)\rangle.
\]

\section*{Complex Conjugates and Modulus}
The \textbf{complex conjugate} reflects a complex number across the real axis:
\[
  z=a+bi\longrightarrow\overline z=a-bi.
\]
In pair form,
\[
  \langle a,b\rangle\longrightarrow\langle a,-b\rangle.
\]
Multiplying a complex number by its conjugate produces a nonnegative real number:
\[
\begin{aligned}
  z\overline z
  &=(a+bi)(a-bi)\\
  &=a^2-abi+abi-b^2i^2\\
  &=a^2+b^2\\
  &=|z|^2.
\end{aligned}
\]
Consequently,
\[
  \boxed{z\overline z=|z|^2},
  \qquad
  |z|=|\overline z|,
  \qquad
  z+\overline z=2\operatorname{Re}(z),
  \qquad
  z-\overline z=2i\operatorname{Im}(z).
\]

Example: if $z=3-4i$, then
\[
  \overline z=3+4i,
  \qquad |z|=5,
  \qquad z\overline z=25.
\]

\section*{Complex Zeros of Polynomials}
\subsection*{Conjugate Pair Theorem}
If a polynomial has \textbf{real coefficients} and $a+bi$ is a nonreal zero, then its conjugate $a-bi$ is also a zero:
\[
  P(a+bi)=0 \quad\Longrightarrow\quad P(a-bi)=0.
\]
The real-coefficient condition matters. Nonreal zeros of real polynomials occur in conjugate pairs.

The pair produces a real quadratic factor:
\[
\begin{aligned}
  \bigl(x-(a+bi)\bigr)\bigl(x-(a-bi)\bigr)
  &=(x-a)^2+b^2\\
  &=x^2-2ax+(a^2+b^2).
\end{aligned}
\]
For example, if $2+3i$ is a zero of a real polynomial, then $2-3i$ is also a zero, and
\[
  \bigl(x-(2+3i)\bigr)\bigl(x-(2-3i)\bigr)
  =(x-2)^2+9=x^2-4x+13
\]
is a factor.

\subsection*{Fundamental Theorem of Algebra}
Every nonconstant polynomial of degree $n$ has exactly $n$ complex zeros when multiplicity is counted. Equivalently, over $\mathbb C$,
\[
  P(x)=a_n(x-r_1)(x-r_2)\cdots(x-r_n).
\]
A repeated zero is counted more than once. For example, if
\[
  P(x)=(x-3)^2(x^2+4),
\]
then its four zeros, counting multiplicity, are $3,3,2i,-2i$.

Useful consequences:
\begin{itemize}[leftmargin=1.5em,itemsep=0.2em]
  \item A quadratic has two complex zeros, a cubic has three, and a quartic has four, counting multiplicity.
  \item A real polynomial of odd degree must have at least one real zero because its nonreal zeros pair up.
  \item For a real polynomial, factor over the reals using linear factors and irreducible quadratic factors; factor over the complex numbers using only linear factors.
  \item The quadratic formula still applies when $b^2-4ac<0$, using $\sqrt{-m}=i\sqrt m$ for $m>0$.
\end{itemize}

\section*{Reciprocals and Division}
For $z=a+bi\ne0$, multiply by the conjugate to write the reciprocal in $a+bi$ form:
\[
\begin{aligned}
  \frac1z
  &=\frac1{a+bi}\cdot\frac{a-bi}{a-bi}\\
  &=\frac{a-bi}{a^2+b^2}\\
  &=\boxed{\frac{a}{a^2+b^2}-\frac{b}{a^2+b^2}i}.
\end{aligned}
\]
Equivalently,
\[
  \frac1z=\frac{\overline z}{|z|^2}.
\]

For $w=c+di$ and nonzero $z=a+bi$,
\[
\begin{aligned}
  \frac wz
  &=\frac{w\overline z}{z\overline z}
  =\frac{(c+di)(a-bi)}{a^2+b^2}\\
  &=\boxed{\frac{ac+bd}{a^2+b^2}
  +\frac{ad-bc}{a^2+b^2}i}.
\end{aligned}
\]
For example,
\[
  \frac{2+i}{1-2i}
  =\frac{(2+i)(1+2i)}{(1-2i)(1+2i)}
  =\frac{5i}{5}=i.
\]

\clearpage
\section*{30-Minute Practice}
\textbf{Instructions.} Complete all 12 questions in 30 minutes. Show clear work. Give exact values, write final complex-number answers in $a+bi$ form unless another form is requested, and use $0\le\operatorname{Arg}z<2\pi$.

\begin{enumerate}[leftmargin=*,itemsep=1em]
  \item Let $z=\langle2,-3\rangle$ and $w=\langle-1,4\rangle$.
  \begin{enumerate}[label=(\alph*),itemsep=0.3em]
    \item Find $2z+w$ in pair form and in $a+bi$ form.
    \item Find $zw$ using the special pair multiplication, then verify using $a+bi$ notation.
  \end{enumerate}

  \item Use the pair multiplication rule to show both of the following directly:
  \[
    \langle0,1\rangle^2=\langle-1,0\rangle,
    \qquad
    \langle r,0\rangle\langle a,b\rangle=\langle ra,rb\rangle.
  \]
  Explain what each identity means in $a+bi$ notation.

  \item For $z=-\sqrt3+i$, find $|z|$ and $\operatorname{Arg}z$. Then write $z$ in polar form
  \[
    |z|(\cos\theta+i\sin\theta).
  \]

  \item Let $z=2-3i$ and $w=-1+4i$.
  \begin{enumerate}[label=(\alph*),itemsep=0.3em]
    \item Find $\overline z$, $|z|$, and $z\overline z$.
    \item Write $1/z$ in $a+bi$ form.
    \item Write $w/z$ in $a+bi$ form.
  \end{enumerate}

  \item Find the real and imaginary parts of
  \[
    \frac{(3-2i)^2}{1+i}.
  \]

  \item A complex number $z$ satisfies
  \[
    z+\overline z=6,
    \qquad
    z\overline z=25.
  \]
  Find all possible values of $z$ and the principal argument of each.

  \item Solve $x^2-6x+13=0$. Factor the quadratic completely over $\mathbb C$.

  \item A monic quadratic polynomial with real coefficients has $-2+3i$ as a zero. Find the other zero and write the polynomial in standard form.

  \item Factor completely over $\mathbb C$ and list all zeros:
  \[
    x^3-5x^2+9x-5.
  \]

  \item Factor completely over the reals and then over $\mathbb C$:
  \[
    x^4-8x^3+26x^2-48x+45.
  \]
  State every zero and its multiplicity.

  \item Find the real quartic polynomial $P(x)$ of least degree such that $2$ is a zero of multiplicity $2$, $-1+2i$ is a zero, and $P(0)=40$. Write $P(x)$ in factored and standard forms.

  \item The polynomial
  \[
    P(x)=x^3+ax^2+bx+10
  \]
  has real coefficients, and $1+2i$ is a zero. Find $a$, $b$, and all three zeros of $P$.
\end{enumerate}

\end{document}
